% 第2章 2.2.3：有限观测能否唯一确定候选空间中的函数。
% (a) 仿射函数 f=1+bt，b=-1,0,1,2，只读 f(0)=1 时全部一致；
% (b) 再读 f(1)=2，唯一确定 b=1；
% (c) 二次多项式 f=1+t+s t(t-1)，s=-2,0,2，读 f(0)=1、f(1)=2 仍全部一致。
\begin{tikzpicture}[font=\figurefont\footnotesize]
\pgfplotsset{every axis/.style={width=40mm,height=36mm,axis lines=left,
  x axis line style={draw=FigureStroke,line width=.6pt},y axis line style={draw=FigureStroke,line width=.6pt},
  tick style={draw=FigureStroke},tick label style={font=\figurefont\scriptsize,text=FigureInk},
  xlabel={$t$},xlabel style={font=\figurefont\scriptsize,at={(axis description cs:1,0)},anchor=north west,yshift=1mm},
  xmin=0,xmax=1.05,ymin=-0.2,ymax=3.2,xtick={0,1},ytick={0,1,2,3},
  ymajorgrids,grid style={draw=FigureGrid,line width=.4pt},
  title style={font=\figurefont\footnotesize},samples=60,domain=0:1}}
\begin{axis}[name=A,title={(a) 只读 $f(0)$}]
\addplot[draw=FigureBlue,line width=1.1pt] {1-x};
\addplot[draw=FigureBlue,line width=1.1pt] {1};
\addplot[draw=FigureBlue,line width=1.1pt] {1+x};
\addplot[draw=FigureBlue,line width=1.1pt] {1+2*x};
\addplot[only marks,mark=*,mark size=1.6pt,FigureInk] coordinates {(0,1)};
\end{axis}
\begin{axis}[name=B,at={(A.east)},anchor=west,xshift=8mm,title={(b) 读 $f(0)$ 与 $f(1)$}]
\addplot[draw=FigureBlueFill,line width=1.1pt] {1-x};
\addplot[draw=FigureBlueFill,line width=1.1pt] {1};
\addplot[draw=FigureBlueFill,line width=1.1pt] {1+2*x};
\addplot[draw=FigureBlue,line width=1.3pt] {1+x};
\addplot[only marks,mark=*,mark size=1.6pt,FigureInk] coordinates {(0,1) (1,2)};
\end{axis}
\begin{axis}[at={(B.east)},anchor=west,xshift=8mm,title={(c) 二次候选，两次读数}]
\addplot[draw=FigureGreen,line width=1.1pt,dashed] {1+x-2*x*(x-1)};
\addplot[draw=FigureBlue,line width=1.3pt] {1+x};
\addplot[draw=FigureGreen,line width=1.1pt,dashed] {1+x+2*x*(x-1)};
\addplot[only marks,mark=*,mark size=1.6pt,FigureInk] coordinates {(0,1) (1,2)};
\end{axis}
\end{tikzpicture}
